\subsection{Diseño de la sesión de pruebas}
\label{sec:diseno-pruebas}

La sesión de pruebas con usuarios se llevó a cabo el 18 de noviembre de 2026. % TODO-CONFIRMAR: escuela (puede anonimizarse), duración de la sesión y quién aplicó las encuestas.
Su objetivo fue valorar la aceptación de la aplicación por parte de los alumnos y la pertinencia pedagógica percibida por los docentes, como complemento del protocolo de validación descrito en la Sección~\ref{sec:sesion-validacion-legibilidad}.

\subsubsection{Participantes}
Participaron 17 alumnos de primer grado de primaria y 2 docentes. Uno de los docentes imparte primer grado y cuenta con 11 a 20 años de experiencia; el otro reporta impartir de primero a sexto grado y cuenta con más de 20 años de experiencia. Ninguno de los alumnos registró los campos opcionales de grado y sexo, por lo que no se reporta su distribución. La muestra se obtuvo por conveniencia y no pretende ser representativa de la población de primer grado.

% TODO-CONFIRMAR: indicar si se recabó el consentimiento de madres, padres o tutores legales y el asentimiento de los alumnos.
Ambos instrumentos fueron anónimos: no se registraron nombres ni datos que permitan identificar a los participantes. Para su análisis, cada encuesta se identificó con un código (A01--A17 para alumnos y D01--D02 para docentes).

\subsubsection{Procedimiento}
Durante la sesión, los alumnos utilizaron la vista del alumno para realizar ejercicios de escritura en la pantalla del dispositivo. % TODO-CONFIRMAR: dispositivo utilizado (tableta, teléfono o computadora) y ejercicios aplicados.
Al finalizar, respondieron la encuesta en papel; el docente leyó cada pregunta en voz alta, ya que en primer grado la lectura autónoma aún está en proceso de consolidación. A los docentes se les presentó el proyecto y las vistas del alumno y del tutor, incluidas la calificación sugerida y la vista de progreso, y después respondieron su propia encuesta.

\subsubsection{Instrumentos}
\textbf{Encuesta para alumnos.} Consta de siete preguntas cerradas con una escala pictórica de tres caritas (\textit{Sí}, \textit{Más o menos} y \textit{No}), formato habitual para recoger opiniones de niños pequeños \cite{read2006funtoolkit}. Las preguntas abordan el agrado general, la comprensión de la actividad, la facilidad de escribir en la pantalla, la lectura de las letras, la utilidad de escuchar las instrucciones y de los dibujos, y la disposición a volver a usar la aplicación. Incluye además dos preguntas de opción única, sobre el medio en que prefieren escribir y el elemento que más les gustó, y una pregunta abierta sobre lo más difícil, que podían responder por escrito o con un dibujo.

\textbf{Encuesta para docentes.} Registra el grado que imparte el docente, sus años de experiencia y la frecuencia con que utiliza tecnología en clase. Contiene doce afirmaciones valoradas en una escala tipo Likert de cinco puntos \cite{likert1932technique}, donde 1 significa \textit{totalmente en desacuerdo} y 5 \textit{totalmente de acuerdo}. Las afirmaciones cubren la claridad de la explicación, la relevancia del problema, la adecuación de los ejercicios al nivel y a los métodos de enseñanza, la facilidad de uso para un niño de 6 a 7 años, el tamaño de textos e íconos, el lienzo, el botón «Escuchar ejercicio», la claridad de la calificación sugerida, la confirmación docente, la vista de progreso y la disposición a usar la aplicación. Se completa con preguntas sobre la frecuencia de uso prevista, las dificultades para usarla en la escuela y una pregunta abierta de mejoras.

\subsubsection{Criterios de registro}
Las respuestas se transcribieron de las hojas originales. Cuando un alumno marcó dos opciones con trazos de intensidad similar en la misma pregunta, la respuesta se registró como \textit{dos marcas} y no se asignó a ninguna opción. Cuando una de las marcas era claramente más tenue, por haber sido borrada o apenas iniciada, se tomó la marca principal. Los porcentajes se calculan sobre el total de encuestas (17) y se presentan redondeados a enteros.
